In Table~\ref{tab-results} we show some experimental results for our
realisation algorithm.

We have performed experiments for sets of scenarios with three
different kinds of sizes, shown in the first column.
%: less than 10, between 10 and 20, and greater than 20.
For each size we considered three categories
of scenarios according to their length, shown in the top row.
%: those whose length is at most 10, between 10 and 15, and
%between 15 and 20.
%Our experiment included 5 scenarios in each group, i.e., within each set
%size and scenario size.
For each kind of size and length we tried 5 different examples.
%experiments.
Table~\ref{tab-results} shows the minimum, maximum and
average number of alternatives of the expressions generated for that
set of examples, as well as the approximate running time of our entire program.

We compared these results with those from a program that performs
a breadth-first search for an optimal solution (i.e., one with the minimum
number of alternatives and as well-balanced as possible): in some of
the cells those results are
shown to the right of the vertical line.
We were unable to do this for all the examples:
%for sets of size larger than 20 the search can take more than a day.
for larger data the complete search becomes prohibitively long.

In two examples our algorithm yields 8 and 11 alternatives as opposed to the
optimal 6 and 10 (see the second row of column one and column two).

%These experiments show that our strategy for quickly finding an
%expression that realises a set of scenarios is effective.

%These experiments illustrate the effectiveness of our strategy for
%finding a scenario expression that realises a set of timed scenarios.

Our experimental results confirm that our approach is very
effective in identifying the divisions that lead to satisfactory
realisations of sets of timed scenarios.

\begin{table}[t]
\begin{center}
\begin{tabular}{ |c|c|c|c| }
  \hline
  $|\mathcal{S}|$ & $5\leq|\SeqSC|\leq10$ & $10<|\SeqSC|\leq 15$&
  $15<|\SeqSC|$\\
  \hline
  $\leq 10$ &
   \begin{tabular}{r|r}
     2, 0.2 s  & 2, \hspace{1em}10.0 s\\
     6, 0.4 s  & 6, \hspace{1em}51.4 s\\
     4.2, 0.3 s & 4.2, \hspace{1em}23.6 s
   \end{tabular}
   &
   \begin{tabular}{r r | r r}
     2, & 0.2 s &   2, & 0.5 s\\
     8, & 1.8 s &   8, & 278.2 s\\
     5.2, & 1.1 s & 5.2, & 65.4 s
   \end{tabular}
     &
   \begin{tabular}{r r | r r}
       3, & 3 s   &   3, & 284.8 s\\
       7,   & 5 s   &   7, & 52.3 s\\
     4.8, & 3.8 s & 4.8, & 171.1 s
     \end{tabular}\\
   \hline
   $\leq 20$ &
   \begin{tabular}{r r | r r}
       3, & 0.4 s &   3, & 1287.2 s\\
       8, & 0.8 s &   6, & 1477.0 s\\
     4.8, & 0.5 s & 4.4, & 4703.8 s
   \end{tabular}
   &
   \begin{tabular}{r r | r r}
       3, & 2.0 s & 3,  & 7401.7 s\\
      11, & 2.9 s & 10, & 147597.7 s\\
     5.4, & 2.2 s & 5.2, & 32438.6 s
   \end{tabular}
   &
   \begin{tabular}{r r | r r}
       2, & 2.6 s & 2, & 17866.5 s\\
      10, & 10.8 s & 10,& 251865.8 s\\
     6.8, & 15.1 s & 6.8, & 92812 s
   \end{tabular}\\
   \hline
   $> 20$ &
   \begin{tabular}{r r}
       3, & 0.9 s\\
      10, & 1.6 s\\
     5.8, & 1.2 s
   \end{tabular} &
   \begin{tabular}{r r}
       3, & 2 s\\
      12, & 4.8 s\\
     6.6, & 3.2 s
   \end{tabular}
   &
   \begin{tabular}{r r}
        3, & 4.4 s\\
       17, & 642.9 s\\
     8.25, & 247.6 s
   \end{tabular} \\
   \hline
\end{tabular}
%\vspace{4mm}
\caption{Experimental results}
\label{tab-results}
%\vspace{-5mm}
\end{center}
\end{table}
